\section{A Model of Priced Equivocation and the Accountability Race}
\label{sec:design}

We now build the model. Table~\ref{tab:notation} summarizes notation, and Table~\ref{tab:results}
previews the formal results, which we prove in this section and \S\ref{sec:reconf}.

\begin{table}[t]
\centering
\caption{Summary of formal results. The single conditional inequality $\Tsettle<\Tacc$ (gate G1)
governs all of them: closing it yields immunity, leaving it open sustains the attack.}
\label{tab:results}
\small
\begin{tabularx}{\linewidth}{@{}l X l@{}}
\toprule
Result & Statement & Where \\
\midrule
Lemma~\ref{lem:intersection} & Two conflicting commit certificates force $\ge F{+}1$ equivocators
(quorum intersection) & \S\ref{sec:floor} \\
Prop.~\ref{prop:immunity} & Finality-gating cross-shard acceptance forces $\pwin{=}0$, hence
$\mathbb{E}[\Pi]\le0$ (immunity by construction) & \S\ref{sec:race} \\
Prop.~\ref{prop:reconf-safe} & A handoff meeting conditions (i)--(iv) keeps $\pwin{=}0$ for old-epoch
certified artifacts after rotation & \S\ref{sec:reconf} \\
Prop.~\ref{prop:reconf-residual} & Ungated carryover makes $\mathbb{E}[\Pi(a)]$ rise toward the
epoch boundary (residual vulnerability) & \S\ref{sec:reconf} \\
Cex.~\ref{cex:persistent}--\ref{cex:rolling} & Four boundaries where the attack fails: persistent
slashability, gated carryover, cross-shard freeze, rolling reconfiguration & \S\ref{sec:reconf} \\
\bottomrule
\end{tabularx}
\end{table}

\begin{table}[t]
\centering
\caption{Notation for the sub-threshold bribery and accountability-race model.}
\label{tab:notation}
\small
\begin{tabularx}{\linewidth}{@{}lX@{}}
\toprule
Symbol & Meaning \\
\midrule
$N=3F+1$ & committee size; $F$ tolerated faults \\
$Q=2F+1$ & commit quorum size \\
$k=F+1$ & equivocators needed for two conflicting certificates \\
$\varphi$ & view-split: fraction of honest votes on one branch \\
$s_i,\,R_i$ & validator $i$'s bonded stake; continuation value \\
$q$ & Pr[fraud proof enforced before stake unslashable] \\
$b_i,\,B$ & bribe to validator $i$; total bribe $\sum_i b_i$ \\
$\Eval$ & value extractable via the conflicting finality \\
$\Tsettle$ & time to irreversibly realize $\Eval$ \\
$\Tacc$ & accountability latency (detect$\rightarrow$slash$\rightarrow$revert) \\
$\pwin$ & $\Pr[\Tsettle < \Tacc]$ \\
$\psucc$ & $\Pr[\text{conflicting finality}\mid k\text{ equivocators}]$ \\
$\Cop$ & operating cost (gas, infrastructure, stake acquisition) \\
$\Pi$ & adversary net profit \\
\bottomrule
\end{tabularx}
\end{table}

\subsection{The Equivocation Floor (Necessary, Not Sufficient)}
\label{sec:floor}

\begin{lemma}[Equivocation floor]\label{lem:intersection}
In a committee of $N=3F+1$ with commit quorum $Q=2F+1$, any two conflicting commit certificates
share at least $F+1$ signers, each of whom must have signed both branches. Hence at least
$k=F+1$ validators must equivocate.
\end{lemma}
\begin{proof}
Two certificates are sets $Q_1,Q_2$ of signers with $|Q_1|=|Q_2|=2F+1$ drawn from $N=3F+1$ validators.
By inclusion--exclusion, $|Q_1\cap Q_2|\ge |Q_1|+|Q_2|-N = 2(2F+1)-(3F+1)=F+1$. A validator in
$Q_1\cap Q_2$ signed both (conflicting) values, i.e.\ equivocated. Thus $|\{\text{equivocators}\}|\ge F+1$.
\end{proof}

Lemma~\ref{lem:intersection} gives a \emph{necessary} condition. It is emphatically \emph{not}
sufficient: each certificate still needs $Q-(F{+}1)=F$ additional signatures from honest validators,
which requires that the honest validators be \emph{split} across the two branches. We make this
quantitative. Suppose the adversary buys $k\ge F+1$ equivocators; the remaining $H=N-k$ honest
validators each commit to one branch, voting for branch $A$ with probability $\varphi$ (the
view-split the adversary engineers via leader equivocation, message delay, or transient asynchrony).
Branch $A$ finalizes iff $k+h_A\ge Q$ and branch $B$ iff $k+(H-h_A)\ge Q$, where $h_A$ is the number
of honest $A$-votes. Both hold iff
\begin{equation}
Q-k \;\le\; h_A \;\le\; N-Q,
\label{eq:band}
\end{equation}
where the upper bound is the branch-$B$ quorum constraint $k+(H-h_A)\ge Q$ with $H=N-k$, giving
$h_A\le H-(Q{-}k)=N-Q$ \emph{independently of} $k$. This is a feasible band of width $k-F$ for
$F{+}1\le k\le Q$ (exactly $1$ when $k=F+1$, the single point $h_A=N-Q=F$): each equivocator bought
beyond the floor adds exactly one feasible split. Modeling $h_A\sim\mathrm{Binomial}(H,\varphi)$,
\begin{equation}
\psucc(N,k,\varphi) \;=\; \Pr\!\big[\,Q-k \le h_A \le N-Q\,\big].
\label{eq:psucc}
\end{equation}
Two consequences matter. First, at the floor $k=F+1$ the band collapses to the single point
$h_A=F$: an adversary with \emph{no} control over views must hit an exact even split, so $\psucc$ is
modest (for $N=31$, $\psucc=\ArtifactPsuccFloor$ at $\varphi=\tfrac12$). Second, each extra purchased
equivocator widens the band, and $\psucc$ is maximized at $\varphi=\tfrac12$. The adversary therefore
trades bribe budget against split-control, and the two routes to a high $\psucc$ price it very
differently. By an \emph{uncontrolled} split, reaching $\psucc\approx0.85$ at $N=31$ needs $k\approx
F{+}6$ (about half the committee), which raises the bribe to ${\approx}1.45\times$ the $F{+}1$ floor
(at $k{=}F{+}5$, ${\approx}1.36\times$, $\psucc\approx0.79$);
by \emph{view/leader control}---the leader-occupation lever the threat model grants only with bounded
probability (\S\ref{sec:threat})---a near-even split is forced at the floor bribe. Consequently the
profitability analyses of \S\ref{sec:eval} that fix $\psucc=0.85$ adopt the view-control regime, and we
bracket every such result against the uncontrolled-split floor $\psucc=\ArtifactPsuccFloor$ rather than
treat $0.85$ as a default. We treat
$\psucc$ as a protocol-and-network property, not a per-validator cost, and validate
Eq.~\eqref{eq:psucc} against a Monte-Carlo BFT harness in \S\ref{sec:eval}.

\subsection{Bribe-Cost Model}
\label{sec:cost}

Under Assumption~\ref{ass:rational}, validator $i$ equivocates iff the bribe covers its expected
loss. Equivocation is maximally accountable (Assumption~\ref{ass:accountable}): if the proof is
enforced before $i$'s stake becomes unslashable, $i$ loses bonded stake $s_i$ and forfeits
continuation value $R_i$. Let $q=\Pr[\text{proof enforced before } i\text{'s stake unslashable}]$.
Then the individual participation constraint is
\begin{equation}
b_i \;\gtrsim\; q\,(s_i + R_i),
\label{eq:bribe-individual}
\end{equation}
and, summing over the $k=F+1$ required equivocators (homogeneous case $s_i=s,\,R_i=R$),
\begin{equation}
B \;=\; \sum_{i=1}^{F+1} b_i \;\approx\; (F+1)\,q\,(s + R).
\label{eq:bribe-total}
\end{equation}
Equation~\eqref{eq:bribe-total} departs from analyses that equate security with total stake: the
adversary pays for $F+1$ validators, not the whole committee, and pays a \emph{fraction} $q$ of each
one's stake-at-risk rather than its face value. Crucially, $q$ is not a constant of the universe. It
is determined by the defender, through (i) the accountability latency $\Tacc$ relative to the
unbonding/withdrawal delay, (ii) the validity and propagation of equivocation evidence, and (iii) the
reliability of enforcement. We make the latency dependence explicit in \S\ref{sec:race} and the
unbonding/rotation dependence in \S\ref{sec:reconf}, where $q$ becomes a function $q(a)$ of attack
time within the epoch.

\subsection{Value Model}
\label{sec:value}

Let $\Eval$ be the value the adversary can extract by acting on the fraudulently finalized branch
before the conflict is resolved. $\Eval$ is upper-bounded by (i) the liquidity reachable through the
target shard's cross-shard and bridge interfaces within $\Tsettle$, and (ii) the state that the
conflicting certificate can fork (e.g.\ the balances whose double-spend the two branches encode). We
treat $\Eval$ as an exogenous deployment parameter and sweep it in \S\ref{sec:eval}.

\subsection{The Accountability Race}
\label{sec:race}

Start the clock when the conflicting certificate exists ($t=0$). Two latencies compete:
\begin{align*}
\Tsettle &= \text{(dest.\ receipt processing)} \\
         &\quad + \text{(withdrawal beyond clawback)},\\
\Tacc    &= \text{(observe both certs)} + \text{(propagate proof)} \\
         &\quad + \text{(enforce slashing)} + \text{(revert/quarantine)}.
\end{align*}
The adversary retains $\Eval$ only if it settles before accountability lands. Define
$\pwin=\Pr[\Tsettle<\Tacc]$. Combining with the success probability $\psucc$ that $k$ equivocators
actually yield conflicting finality, the realized value is $\psucc\,\pwin\,\Eval$ in expectation. This
product treats the honest-split event (probability $\psucc$) and the window-open event (probability
$\pwin$) as \emph{independent}; we adopt this as a modeling assumption and note its sign. Both are
driven by the same lever---induced asynchrony or leader control raises $\psucc$ and, by delaying
certificate observation and proof propagation, also raises $\pwin$ (the $\Tacc$-as-master-variable
coupling below)---so their covariance is plausibly nonnegative, making the factored $\psucc\,\pwin$ a
\emph{lower} bound on $\mathbb{E}[\psucc\pwin]$ and the break-even $\Eval^{*}\sim1/(\psucc\pwin)$ a
conservative \emph{upper} bound. Immunity ($\pwin=0$) is unaffected.

\paragraph{Profitability.}
With \emph{guided} bribes (paid regardless of outcome, per Assumption~\ref{ass:rational} and
\cite{karakostas2024bribing}) and operating cost $\Cop$, expected profit is
\begin{equation}
\mathbb{E}[\Pi] \;=\; \psucc\,\pwin\,\Eval \;-\; B \;-\; \Cop.
\label{eq:profit-guided}
\end{equation}
With \emph{effective} (success-conditional) bribes, the cost moves inside the success event but
validators demand a premium $B_{\mathrm{succ}}>B$ to bear slashing risk on failed attempts, while
failed attempts may still incur $B_{\mathrm{fail}}\ge 0$:
\begin{equation}
\begin{aligned}
\mathbb{E}[\Pi] ={}& \psucc\,(\pwin\,\Eval - B_{\mathrm{succ}})\\
&{} - (1-\psucc)\,B_{\mathrm{fail}} - \Cop.
\end{aligned}
\label{eq:profit-effective}
\end{equation}
The attack is worth attempting iff $\mathbb{E}[\Pi]>0$. In the deterministic-win, full-success limit
($\psucc=\pwin=1$), \eqref{eq:profit-guided} reduces to a clean pair of \emph{independent gates}:
\begin{align}
\mathrm{(G1)}\quad & \Tsettle < \Tacc, \label{eq:g1}\\
\mathrm{(G2)}\quad & \Eval > B + \Cop, \label{eq:g2}\\
& B=(F{+}1)\,q\,(s{+}R). \nonumber
\end{align}
Both gates must hold. (G1) is a property of the protocol's pipeline and accountability machinery;
(G2) is an economic property of stake, liquidity, and $q$.

\paragraph{$\Tacc$ as a master variable.}
The two gates are coupled through $\Tacc$. A larger $\Tacc$ (slower or less reliable accountability)
raises $\pwin$ in (G1) \emph{and} lowers $q$ in (G2)---validators who expect to escape enforcement
before unbonding accept smaller bribes. Thus slow accountability harms the defender twice: it widens
the settlement window \emph{and} cheapens corruption. ``Cheap to corrupt'' and ``easy to cash out''
are two symptoms of one underlying quantity. This observation, rather than any single inequality, is
the conceptual core of the paper, and it is why purely economic defenses (bigger committees, more
stake) are not a substitute for closing the window.

\begin{proposition}[Immunity by construction]\label{prop:immunity}
If the protocol admits cross-shard (and bridge) effects predicated on a source-shard certificate
only after that certificate is past the source shard's accountability window---formally, if it
guarantees $\Tsettle \ge \Tacc$ for every such effect---then $\pwin=0$, and by
\eqref{eq:profit-guided} $\mathbb{E}[\Pi]=-B-\Cop\le 0$ for any $B\ge 0,\ \Cop\ge 0$, strictly
negative whenever $B+\Cop>0$ (any attack that pays a bribe or gas). No rational
adversary (Assumption~\ref{ass:rational}) mounts the attack, independent of committee size $N$ or
stake $s$.
\end{proposition}
\begin{proof}
By hypothesis every value-realizing cross-shard effect occurs at or after $\Tacc$, so
$\Pr[\Tsettle<\Tacc]=0$, i.e.\ $\pwin=0$. Substituting into \eqref{eq:profit-guided} gives
$\mathbb{E}[\Pi]=-(B+\Cop)\le 0$, strictly negative whenever the adversary pays any bribe or gas. Under
Assumption~\ref{ass:rational} the adversary declines a strictly negative-expectation action.
\end{proof}

Proposition~\ref{prop:immunity} is a design rule as much as a theorem: it states the exact condition
a sharded system must enforce to be unconditionally safe against this attack, and explains why
enlarging committees or stake (which affect only (G2)) does not help once (G1) is open.

\subsection{Attack Procedure}
\label{sec:procedure}

Algorithm~\ref{alg:attack} states the end-to-end procedure; Figure~\ref{fig:flow} depicts it. The
trustless payment mechanism is a contract that releases a bribe only upon submission of a valid
equivocation proof by a registered validator, so neither party need trust the other
(cf.~\cite{sookitoth2025bribers}). Appendix~\ref{app:contract} sketches the contract; \S\ref{sec:impl}
reports a runnable local model of the same logic.

\begin{figure}[t]
\centering
\includegraphics[width=0.86\linewidth]{fig1_attack_flow.pdf}
\caption{Sub-threshold committee bribery and the accountability race. The adversary retains $V$ only
if cross-shard settlement ($\Tsettle$) precedes end-to-end accountability ($\Tacc$).}
\label{fig:flow}
\end{figure}

\begin{algorithm}[t]
\caption{Sub-threshold committee bribery with accountability race}
\label{alg:attack}
\begin{algorithmic}[1]
\State \textbf{Input:} target committee $C$ ($|C|=3F{+}1$), value target $\Eval$, budget
\State Estimate $\Tacc,\Tsettle,q$ for $C$ and its cross-shard interface
\If{$\Tsettle \ge \Tacc$ \textbf{or} $\Eval \le (F{+}1)q(s{+}R)+\Cop$}
  \State \textbf{abort} \Comment{a gate fails; not profitable (cf.\ Prop.~\ref{prop:immunity})}
\EndIf
\State Deploy \textsc{PayToEquivocate} escrow funded with $B \approx (F{+}1)q(s{+}R)$
\State Recruit $k=F{+}1$ validators of $C$ to register and commit to equivocate
\State Engineer/await a usable view (leader equivocation or transient asynchrony) so honest
       validators split across branches $b_A \ne b_B$
\State Drive $C$ to finalize conflicting certificates $\sigma_A$ on $b_A$ and $\sigma_B$ on $b_B$
\State On branch $b_A$, submit the cross-shard receipt that moves/duplicates $\Eval$
\State Withdraw $\Eval$ on the destination beyond clawback \Comment{completes at $\Tsettle$}
\State Bribed validators submit equivocation proofs to \textsc{PayToEquivocate}; bribes released
\State \textbf{Outcome:} if completed before $\Tacc$, adversary retains $\Eval$; bribed validators
       are slashed
\end{algorithmic}
\end{algorithm}

